%!TEX root = ../main.tex

\chapter{Konzeption der Testinfrastruktur}

\section{Zwei-Phasen-Rollout-Konzept (lokal $\rightarrow$ zentralisiert)}
Die Konzeption folgt einem zweistufigen Vorgehen: zuerst eine belastbare lokale Testinfrastruktur, anschließend die Überführung in eine zentrale Ausführungsumgebung.

\textbf{Leitgedanke:}
Lokale Stabilität und Teamakzeptanz sind Voraussetzung für eine erfolgreiche Zentralisierung.

\subsection{Konzept der lokalen Testumgebung}

\subsubsection{Einordnung des ersten Testkonzepts}
Das erste Testkonzept (Stand: Juni 2026) definiert den Einstieg in eine strukturierte Teststrategie für ATLAS. Es beschreibt eine klare Priorisierung kritischer Komponenten, eine Zuordnung von Unit-, Integrations- und End-to-End-Tests sowie ein zweistufiges Zielbild aus lokaler Umsetzung und späterer CI/CD-Einbindung.

Für den initialen Stack waren unter anderem xUnit, WebApplicationFactory und Moq im Backend, Angular TestBed im Frontend sowie Playwright für End-to-End-Szenarien vorgesehen. Im frühen Entwurf wurden bei E2E-Tests zusätzlich Alternativen betrachtet; im weiteren Projektverlauf wurde die Umsetzung auf Playwright fokussiert.
Die Einordnung dieses Entwicklungsstands basiert auf dem projektspezifischen Testkonzeptentwurf \cite{Pfeil2026Testkonzept}.


Beschrieben werden Setup, Ausführungsabläufe, Testgrenzen und die Übergabe an den Main-Entwickler zur Pipeline-Integration.

Das initiale Konzept priorisiert im Backend zuerst GatewayService, ExperimentService und FileService (Priorität 1), danach ReportService und SearchService (Priorität 2). Im Frontend wurden zunächst die fachkritischen Bereiche Features, Services, Auth und Guards priorisiert; wiederverwendbare Komponenten und Layout-Bausteine wurden nachgelagert betrachtet.

Methodisch wurde eine Testpyramide als Zielbild angesetzt: hoher Anteil schneller Unit-Tests, ergänzende Integrations-Tests und ein kleiner, stabiler Anteil End-to-End-Tests für kritische User Journeys.

\textbf{Zu konkretisieren:}
\begin{itemize}
	\item lokale Ausführung aller relevanten Testarten
	\item einheitliche Projektstruktur und Namenskonventionen
	\item Übergabepunkte für die Integration in CI
\end{itemize}

\subsection{Konzept für die zentrale Testumgebung (Agent/VM)}
Dieser Abschnitt beschreibt das Zielbild einer zentralisierten Testausführung und definiert Entscheidungskriterien für den Umsetzungsgrad.

Im ersten Konzept war die zentrale Ausführung über Azure DevOps mit Agents vorgesehen. Geplant war, Testläufe früh im Pull-Request-Prozess auszuführen und die Ergebnisse als Qualitätsindikator für Merge-Entscheidungen zu nutzen.

\subsubsection{Bewertungsoption: 1:1 gesamte Infrastruktur duplizieren}
Vorteil: hohe Realitätsnähe. Nachteil: höherer Betriebs- und Pflegeaufwand.

\subsubsection{Bewertungsoption: alles außer Datenbank und Schnittstellen duplizieren}
Vorteil: geringere Komplexität. Nachteil: reduzierte Aussagekraft für bestimmte Integrationsszenarien.

\subsubsection{Kritische Reflexion der zentralen Umsetzungsvariante}
Hier wird transparent bewertet, welche Variante unter den Projektbedingungen wirtschaftlich und technisch tragfähig ist.

\section{Mögliche Testherausforderungen}
Mögliche Herausforderungen sind instabile Testdaten, Umgebungsabhängigkeiten, Flaky Tests sowie längere Laufzeiten in zentralen Pipelines.
